\subsection{Network Architectures and Expressivity}

The migrated deep-feedforward-network notes describe neural networks as learned feature maps followed by prediction heads.

\begin{definition} \label{def:sl-feedforward-network} \hlt{Deep Feedforward Network}\\
Let $x\in\mathcal X$, set $h^{(0)}=x$, and fix depth $L$.
\begin{enumerate}[label=(\roman*)]
\setlength{\itemsep}{0pt}
\item \hlt{Layer recursion}: $h^{(k)}=g^{(k)}(W^{(k)}h^{(k-1)}+b^{(k)})$ for $k=1,\ldots,L$.
\item \hlt{Network map}: $f_\theta:\mathcal X\to\mathcal Y$ is $f_\theta(x)=h^{(L)}=(f_L\circ\cdots\circ f_1)(x)$.
\item \hlt{Parameters}: $\theta=\{(W^{(k)},b^{(k)})\}_{k=1}^L$.
\item \hlt{Training objective}: choose $\theta$ to minimize empirical risk $\hat R_n(f_\theta)$ or population risk $R(f_\theta)$.
\end{enumerate}
\end{definition}

\begin{remark} \label{rem:sl-learned-feature-map} \hlt{Learned Feature Maps}\\
Let $\ell$ be the task loss and let $R(f)=\mathbb E[\ell(Y,f(X))]$. Two representation parametrizations are:
\begin{enumerate}[label=(\roman*)]
\setlength{\itemsep}{0pt}
\item \hlt{fixed features}: $f_w(x)=a_w(\phi(x))$, with $\phi$ fixed and only the head parameter $w$ optimized.
\item \hlt{learned features}: $f_{\theta,w}(x)=a_w(\phi_\theta(x))$, with $\theta$ and $w$ optimized jointly.
\end{enumerate}
Learning $\phi_\theta$ adds adaptivity but usually destroys convexity.
\end{remark}

\begin{definition} \label{def:sl-projection-pursuit-ridge-functions} \hlt{Projection Pursuit and Ridge Functions}\\
Let $x\in\R^p$.
\begin{enumerate}[label=(\roman*)]
\setlength{\itemsep}{0pt}
\item \hlt{Projection-pursuit form}: $f(x)=\sum\limits_{m=1}^M g_m(\omega_m^Tx)$ with $\norm{\omega_m}_2=1$.
\item \hlt{Ridge coordinate}: $u_m=\omega_m^Tx$.
\item \hlt{Ridge function}: $g_m(u_m)$ depends on $x$ only through the scalar coordinate $u_m$.
\item \hlt{Learned components}: the directions $\omega_m$ and univariate profiles $g_m$ are estimated from data.
\end{enumerate}
\end{definition}

\begin{definition} \label{def:sl-output-and-hidden-units} \hlt{Output and Hidden Units}\\
Let $T\in\R^K$ denote output scores. Common output parametrizations are:
\begin{enumerate}[label=(\roman*)]
\setlength{\itemsep}{0pt}
\item \hlt{linear}: $f_k(x)=T_k$, for conditional means or real-valued responses.
\item \hlt{sigmoid}: $p(x)=\sigma(T)=\frac{1}{1+\exp(-T)}$, for Bernoulli probabilities.
\item \hlt{softmax}: $p_k(x)=\exp(T_k)/\sum\limits_{\ell=1}^K\exp(T_\ell)$, for $K$-class probabilities.
\end{enumerate}
Hidden units define activations $Z_m=\phi_m(a_m)$ from preactivations $a_m$, commonly using ReLU, leaky ReLU, maxout, $\tanh$, or another nonlinear map.
\end{definition}

\begin{definition} \label{def:sl-single-hidden-layer-network} \hlt{Single Hidden Layer Network}\\
Let $x\in\R^p$, $Z=(Z_1,\ldots,Z_M)^T$, $T=(T_1,\ldots,T_K)^T$. A single-hidden-layer feedforward network is defined by:
\begin{enumerate}[label=(\roman*)]
\setlength{\itemsep}{0pt}
\item $Z_m=\sigma(\alpha_{0m}+\alpha_m^Tx)$ for $m=1,\ldots,M$.
\item $T_k=\beta_{0k}+\beta_k^TZ$ for $k=1,\ldots,K$.
\item $f_k(x)=g_k(T)$, with $g_k(T)=T_k$ for regression and $g_k(T)=\exp(T_k)/\sum\limits_{\ell=1}^K\exp(T_\ell)$ for $K$-class classification.
\end{enumerate}
In the classification case, $f(x)\in\Delta^{K-1}$.
\end{definition}

\begin{figure}[H]
\centering
\begin{tikzpicture}[
    >=Stealth,
    font=\scriptsize,
    line cap=round,
    line join=round,
    block/.style={draw=gray!65, rounded corners=2pt, align=center,
        minimum height=0.72cm, inner xsep=6pt, inner ysep=2pt},
    input/.style={block, draw=cyan!65!black, fill=cyan!8, minimum width=1.65cm},
    hidden/.style={block, draw=blue!55!black, fill=blue!8, minimum width=2.85cm},
    score/.style={block, draw=orange!70!black, fill=orange!10, minimum width=2.65cm},
    flow/.style={->, draw=black!75, line width=0.6pt}
]
\node[input] (x) at (0,0) {inputs\\$x=(X_1,\ldots,X_p)$};
\path (x.east) ++(0.7,0) node[hidden, anchor=west] (z)
  {hidden features\\$Z_m=\sigma(\alpha_{0m}+\alpha_m^Tx)$};
\path (z.east) ++(0.7,0) node[score, anchor=west] (t)
  {output scores\\$T_k=\beta_{0k}+\beta_k^TZ$};

\draw[flow] (x) -- (z);
\draw[flow] (z) -- (t);
\end{tikzpicture}

\caption{Single-hidden-layer feedforward network.}
\label{fig:sl-single-hidden-layer-network}
\end{figure}

\begin{remark} \label{rem:sl-neural-networks-as-ridge-expansions} \hlt{Neural Networks as Ridge Expansions}\\
For scalar output, a one-hidden-layer network has $f(x)=\beta_0+\sum\limits_{m=1}^M\beta_m\sigma(\alpha_{0m}+\alpha_m^Tx)$.\\
If $\alpha_m\ne0$ and $\omega_m=\frac{\alpha_m}{\norm{\alpha_m}_2}$, the $m$th term is $\beta_m\sigma(\alpha_{0m}+\norm{\alpha_m}_2\omega_m^Tx)$, a ridge term in coordinate $\omega_m^Tx$ with learned \hlt{direction} $\omega_m$, \hlt{scale} $\norm{\alpha_m}_2$, \hlt{threshold} $\alpha_{0m}$, and \hlt{output coefficient} $\beta_m$.\\
Thus a shallow scalar network is a projection-pursuit expansion with parametrized ridge profiles.
\end{remark}

\begin{theorem} \label{thm:sl-universal-approximation-form} \hlt{Universal Approximation Theorem Form}\\
For compact $K\subset\R^d$, $h\in C(K)$, non-polynomial activation $\sigma:\R\to\R$, $\forall \varepsilon>0$, $\exists$ $M\in\mathbb{N}$, $a_m,b_m\in\R$, $w_m\in\R^d$ s.t. one-hidden-layer network $f_M(x)=\sum\limits_{m=1}^M a_m\sigma(w_m^\top x+b_m)$ satisfies $\norm{f_M-h}_\infty<\varepsilon$ on $K$.
\begin{proof} Let $\mathcal{A}_\sigma\defeq\operatorname{span}\{\sigma(w^\top x+b):w\in\R^d,\ b\in\R\}\subset C(K)$.\\
The classical universal approximation theorem states that $\overline{\mathcal{A}_\sigma}^{\norm{\cdot}_\infty}=C(K)$ for suitable non-polynomial $\sigma$.\\
Since $h\in C(K)$ and $\varepsilon>0$, there is $f_M\in\mathcal{A}_\sigma$ with $\norm{f_M-h}_\infty<\varepsilon$.\\
Thus $f_M(x)=\sum\limits_{m=1}^M a_m\sigma(w_m^\top x+b_m)$ has one hidden layer and linear output.\\
\hlt{Structural Commentary}: Non-polynomial activations make ridge spans dense in $C(K)$.\\
\hlt{Key Insight}: Width supplies finite linear combinations of nonlinear ridge functions.
\end{proof}
\end{theorem}

\begin{remark} \label{rem:sl-depth-expressivity} \hlt{Depth and Expressivity}\\
Let $N_L(\varepsilon,h)$ be the minimum width of a depth-$L$ network with $\norm{f-h}_\infty\le\varepsilon$ on a compact domain. Universal approximation gives $N_L(\varepsilon,h)<\infty$, not trainability; depth expressivity is $N_L(\varepsilon,h)\ll N_{L'}(\varepsilon,h)$ for some $L>L'$.
\end{remark}
